\section{The physical packet and the matched cutoff shell}
\label{sec:tpc32-physical-shell}

This section fixes the physical object before any determinant-frequency
transform is taken.  The order is important: the two mixed channels and
the both-ultra channel come from one difference of quadratic prefixes and
must first be recombined.  The interface is inherited from
TPC-25 and TPC-29--TPC-31
\citep{WangTPC25,WangTPC29,WangTPC30,WangTPC31}.

\subsection{Primitive opened rows}

Fix \(h\in\Z\setminus\{0\}\), and put
\begin{equation}
 H=\operatorname{rad}|h|.
 \label{eq:tpc32-H}
\end{equation}
Let \(X\) tend to infinity and fix \(0<\delta<1/2\).  In one opened
dyadic packet the scales satisfy
\begin{equation}
 R=X^{1/2-\delta+o(1)},\qquad
 Q=LD,\qquad QJ\asymp_{\mathscr D}X,
 \label{eq:tpc32-physical-scales}
\end{equation}
with
\begin{equation}
 D\le R^{1/2+o(1)},\qquad
 L>\max\{2R,2|h|\}.
 \label{eq:tpc32-row-separation}
\end{equation}
Equivalently, when exponent notation is useful, we write
\begin{equation}
 Q=X^{\beta+o(1)},\qquad
 J=X^{1-\beta+o(1)},\qquad 0<\beta<1.
 \label{eq:tpc32-beta-scales}
\end{equation}
Here \(\mathscr D\) denotes all fixed support, smoothness, residue, and
positive-margin data.

Rows are indexed by
\begin{equation}
 \alpha=(\ell_\alpha,d_\alpha),\qquad
 m_\alpha=\ell_\alpha d_\alpha\asymp_{\mathscr D}Q,
 \label{eq:tpc32-row-index}
\end{equation}
where \(\ell_\alpha\in[L/2,2L]\) is prime and
\(d_\alpha\in[D,2D]\cap[1,R^{1/2+o(1)}]\) is squarefree with
\((d_\alpha,H)=1\).  The size separation makes \(m_\alpha\) determine
\(\alpha\) uniquely.  The actual one-row coefficients have the
prime--M\"obius provenance
\begin{equation}
 \gamma_\alpha^{(i)}
 =\mu(d_\alpha)(\log\ell_\alpha)
  \omega_D^{(i)}(d_\alpha)
  \psi_L^{(i)}(\ell_\alpha/L)\zeta_\alpha^{(i)},
 \qquad |\zeta_\alpha^{(i)}|\ll_{\mathscr D}1,
 \quad i=1,2.
 \label{eq:tpc32-row-provenance}
\end{equation}
The smooth dyadic families are fixed and bounded.  In particular
\begin{align}
 \sum_\alpha|\gamma_\alpha^{(i)}|
 &\ll_{\mathscr D}Q,
 \label{eq:tpc32-row-l1}\\
 \sum_\alpha|\gamma_\alpha^{(i)}|^2
 &\ll_{\mathscr D}Q\log(2L),
 \qquad i=1,2.
 \label{eq:tpc32-row-l2}
\end{align}

For an orbit point \(j\), set
\begin{equation}
 N_\alpha(j)=m_\alpha j+h.
 \label{eq:tpc32-target}
\end{equation}
Every contributing row and orbit point obeys
\begin{equation}
 (m_\alpha,H)=(j,H)=1,
 \qquad
 c_{\mathscr D}X\le N_\alpha(j)\le C_{\mathscr D}X,
 \label{eq:tpc32-primitive-target-support}
\end{equation}
and the joint \(j\)-support contains \(O_{\mathscr D}(J)\) integers.
Consequently every divisor \(w\mid N_\alpha(j)\) satisfies
\begin{equation}
 (w,m_\alpha H)=1.
 \label{eq:tpc32-target-divisor-primitivity}
\end{equation}

The row-pair mask and the remaining physical weights are retained in
the single joint multiplier
\begin{equation}
 \mathfrak A_{\alpha,\gamma}(j)
 :=\mathfrak m(\alpha,\gamma)
   \Xi_{\alpha,\gamma}(j)
   \cW_{\alpha,\gamma}(j/J).
 \label{eq:tpc32-joint-multiplier}
\end{equation}
We use only
\begin{equation}
 |\mathfrak A_{\alpha,\gamma}(j)|\ll_{\mathscr D}1,
 \qquad
 \mathfrak A_{\alpha,\gamma}(j)=0
 \quad\text{if }m_\alpha=m_\gamma,
 \label{eq:tpc32-joint-multiplier-bound}
\end{equation}
and that \(\mathfrak m\) is independent of the divisor variables.  The
fixed-residue and smooth parts of
\eqref{eq:tpc32-joint-multiplier} possess the finite Fourier and Mellin
separations recorded in TPC-25.  No separability of the generic
row-pair mask is assumed here.

For a kernel \(\cK_{\alpha,\gamma}(j)\) and an event
\(\cE(\alpha,\gamma,j)\), define
\begin{align}
 \cS[\cK;\cE]
 :={}&\sum_{\alpha,\gamma}
 \gamma_\alpha^{(1)}\gamma_\gamma^{(2)}
 \sum_j\mathfrak A_{\alpha,\gamma}(j)
 \cK_{\alpha,\gamma}(j)
 \one_{\cE(\alpha,\gamma,j)}.
 \label{eq:tpc32-physical-functional}
\end{align}
The natural absolute scale is
\begin{equation}
 JQ^2\asymp XQ.
 \label{eq:tpc32-natural-scale}
\end{equation}

\subsection{Divisor coefficients and exact prefixes}

Put
\begin{align}
 a(w)&=-\mu(w)\log w,
 \label{eq:tpc32-a-coefficient}\\
 \lambda'_R(w)
 &:={}
 \begin{cases}
 \displaystyle
 w\sum_{b\le R/w}\frac{\mu(wb)\mu(b)}{\varphi(wb)},&w\le R,\\[7pt]
 0,&w>R,
 \end{cases}
 \label{eq:tpc32-lambda-coefficient}\\
 b_R(w)&=a(w)-\lambda'_R(w).
 \label{eq:tpc32-b-coefficient}
\end{align}
All three coefficients vanish off the squarefree integers.  On
squarefree support,
\begin{align}
 b_R(w)&=-\mu(w)w_R(w),
 \label{eq:tpc32-positive-coefficient-form}\\
 w_R(w)
 &:={}
 \log w+\one_{w\le R}\frac{w}{\varphi(w)}
 \sum_{\substack{b\le R/w\\(b,w)=1}}
 \frac{\mu(b)^2}{\varphi(b)}>0.
 \label{eq:tpc32-wR}
\end{align}

Choose integers
\begin{equation}
 R\le T<U_0,
 \qquad U_0\asymp_{\mathscr D}X,
 \label{eq:tpc32-cutoff-range}
\end{equation}
with \(U_0\) larger than every target on
\eqref{eq:tpc32-primitive-target-support}.  Define the raw prefix
\begin{equation}
 \mathsf A_{m,Y}(j)
 :=\sum_{\substack{w\le Y\\w\mid mj+h}}b_R(w),
 \qquad R\le Y\le U_0,
 \label{eq:tpc32-raw-prefix}
\end{equation}
and the ultra increment
\begin{equation}
 C_m(j)
 :=\sum_{\substack{T<w\le U_0\\w\mid mj+h}}a(w).
 \label{eq:tpc32-ultra-increment}
\end{equation}
Since \(T\ge R\), one has \(b_R(w)=a(w)\) throughout the new shell.
Thus, pointwise in \(m,j\),
\begin{equation}
 \boxed{
 \mathsf A_{m,U_0}(j)
 =\mathsf A_{m,T}(j)+C_m(j).}
 \label{eq:tpc32-exact-prefix-increment}
\end{equation}
The elementary divisor estimate gives, uniformly on the physical
support,
\begin{equation}
 |\mathsf A_{m,T}(j)|+|\mathsf A_{m,U_0}(j)|+|C_m(j)|
 \ll_{\eps,\mathscr D}X^\eps.
 \label{eq:tpc32-prefix-envelope}
\end{equation}

The actual row coefficient and the divisor coefficient fuse on each
leg without an extra hypothesis.

\begin{lemma}[Prime--M\"obius sign fusion]
\label{lem:tpc32-sign-fusion}
If \(w\mid N_\alpha(j)\) and \(b_R(w)\ne0\), then
\begin{equation}
 \boxed{
 \mu(d_\alpha)b_R(w)
 =-\mu(d_\alpha w)w_R(w).}
 \label{eq:tpc32-sign-fusion}
\end{equation}
\end{lemma}

\begin{proof}
The support of \(b_R\) makes \(w\) squarefree, while
\(d_\alpha\) is squarefree by construction.  Since
\(d_\alpha\mid m_\alpha\), target primitivity
\eqref{eq:tpc32-target-divisor-primitivity} gives
\((d_\alpha,w)=1\).  Hence
\(\mu(d_\alpha)\mu(w)=\mu(d_\alpha w)\), and
\eqref{eq:tpc32-positive-coefficient-form} proves the claim.
\end{proof}

\begin{remark}[A fusion which is not available]
The M\"obius factor introduced later by exact-content inversion may
share prime factors with \(w\).  It therefore cannot be fused with
\(\mu(w)\) unless an additional coprimality condition has first been
proved.  No such fusion is used below.
\end{remark}

\subsection{Three raw channels as one rank--two shell}

For two rows \(m=m_\alpha\) and \(n=m_\gamma\), define
\begin{align}
 \cK^{(1)}_{\alpha,\gamma}(j)
 &=\mathsf A_{m,T}(j)C_n(j),
 \notag\\
 \cK^{(2)}_{\alpha,\gamma}(j)
 &=C_m(j)\mathsf A_{n,T}(j),
 \label{eq:tpc32-three-raw-kernels}\\
 \cK^{(3)}_{\alpha,\gamma}(j)
 &=C_m(j)C_n(j).
 \notag
\end{align}
Their matched sum is
\begin{equation}
 \cK^{\mathrm{sh}}_{\alpha,\gamma}(j)
 :=\sum_{\nu=1}^3\cK^{(\nu)}_{\alpha,\gamma}(j).
 \label{eq:tpc32-matched-shell-definition}
\end{equation}

\begin{proposition}[Exact cutoff shell]
\label{prop:tpc32-exact-cutoff-shell}
Pointwise in every physical row pair and orbit point,
\begin{equation}
 \boxed{
 \cK^{\mathrm{sh}}_{\alpha,\gamma}(j)
 =\mathsf A_{m,U_0}(j)\mathsf A_{n,U_0}(j)
  -\mathsf A_{m,T}(j)\mathsf A_{n,T}(j).}
 \label{eq:tpc32-exact-cutoff-shell}
\end{equation}
Equivalently, after opening both divisor sums,
\begin{align}
 \cK^{\mathrm{sh}}_{\alpha,\gamma}(j)
 ={}&\sum_{u,v\ge1}b_R(u)b_R(v)
 \Theta_{T,U_0}(u,v)
 \one_{u\mid N_\alpha(j)}\one_{v\mid N_\gamma(j)},
 \label{eq:tpc32-opened-matched-shell}\\
 \Theta_{T,U_0}(u,v)
 :={}&\one_{u\le U_0}\one_{v\le U_0}
       -\one_{u\le T}\one_{v\le T}.
 \label{eq:tpc32-cutoff-tensor}
\end{align}
In particular \(\Theta_{T,U_0}\), viewed as a matrix in \(u,v\), has
algebraic rank at most two.
\end{proposition}

\begin{proof}
Insert \eqref{eq:tpc32-exact-prefix-increment} on both rows and expand
the difference of products.  This gives the three terms in
\eqref{eq:tpc32-three-raw-kernels}.  Opening the prefixes proves
\eqref{eq:tpc32-opened-matched-shell}.  Finally, with
\(p_Y(u)=\one_{u\le Y}\),
\[
 \Theta_{T,U_0}=p_{U_0}\otimes p_{U_0}-p_T\otimes p_T,
\]
a difference of two rank-one tensors.
\end{proof}

There is a useful symmetric polarization which keeps the same rank
bound.  Put
\begin{equation}
 d_{T,U_0}=p_{U_0}-p_T,
 \qquad
 s_{T,U_0}=\frac{p_{U_0}+p_T}{2}.
 \label{eq:tpc32-cutoff-polarization-vectors}
\end{equation}
Then
\begin{equation}
 \boxed{
 \Theta_{T,U_0}
 =d_{T,U_0}\otimes s_{T,U_0}
  +s_{T,U_0}\otimes d_{T,U_0}.}
 \label{eq:tpc32-cutoff-polarization}
\end{equation}
At the level of the physical divisor sums, define
\begin{equation}
 \mathsf H_{m;T,U_0}(j)
 :=\frac{\mathsf A_{m,U_0}(j)+\mathsf A_{m,T}(j)}2
 =\mathsf A_{m,T}(j)+\frac12C_m(j).
 \label{eq:tpc32-half-prefix}
\end{equation}
Then the exact polarized shell is
\begin{equation}
 \boxed{
 \cK^{\mathrm{sh}}_{\alpha,\gamma}(j)
 =C_m(j)\mathsf H_{n;T,U_0}(j)
  +\mathsf H_{m;T,U_0}(j)C_n(j).}
 \label{eq:tpc32-physical-polarization}
\end{equation}
Thus three rectangles become two polarized tensors before estimation.
This is an algebraic rank statement.  It does not make the masked row
form positive semidefinite and does not authorize deletion of either
polarized term.

\subsection{The complete calibrated difference and its drift boundary}

For \(m=\ell d\), let
\begin{equation}
 \delta_H(m)
 =\frac{H}{\varphi(H)}
  \frac{d}{\varphi(d)(\ell-1)}
 \ll_{\eps,h}\frac{X^\eps}{L},
 \label{eq:tpc32-row-drift}
\end{equation}
and define the calibrated prefix
\begin{equation}
 P_{m,Y}(j)=\mathsf A_{m,Y}(j)-\delta_H(m).
 \label{eq:tpc32-calibrated-prefix}
\end{equation}

\begin{proposition}[Exact calibrated cutoff difference]
\label{prop:tpc32-calibrated-cutoff}
For every physical \(m,n,j\),
\begin{align}
 &P_{m,U_0}(j)P_{n,U_0}(j)
  -P_{m,T}(j)P_{n,T}(j)
 \notag\\
 &\quad={}
 \cK^{\mathrm{sh}}_{\alpha,\gamma}(j)
 -\delta_H(m)C_n(j)-\delta_H(n)C_m(j).
 \label{eq:tpc32-calibrated-cutoff}
\end{align}
There is no two-drift term: it cancels between the two cutoff levels.
\end{proposition}

\begin{proof}
Use \(P_{m,U_0}=P_{m,T}+C_m\), which follows from
\eqref{eq:tpc32-exact-prefix-increment}, and expand.  Alternatively,
insert \eqref{eq:tpc32-calibrated-prefix} into the difference in
\eqref{eq:tpc32-exact-cutoff-shell}; the two copies of
\(\delta_H(m)\delta_H(n)\) cancel.
\end{proof}

Let \(\cD^{\mathrm{dr}}(\cE)\) be the contribution of the final two
terms in \eqref{eq:tpc32-calibrated-cutoff} to
\eqref{eq:tpc32-physical-functional}, with any additional event
\(\cE\).  Since an indicator or a unit-modulus determinant phase can
only decrease its absolute majorant, the following estimate is uniform
over every restriction used later.

\begin{proposition}[Uniform drift boundary]
\label{prop:tpc32-drift-boundary}
For every fixed \(\eps>0\),
\begin{equation}
 \boxed{
 |\cD^{\mathrm{dr}}(\cE)|
 \ll_{\eps,h,\mathscr D}
 X^\eps\frac{XQ}{L}.}
 \label{eq:tpc32-drift-boundary}
\end{equation}
Consequently, if \(L\ge X^{\lambda_0}\) for fixed
\(\lambda_0>0\), the drift legs have a fixed-power saving relative to
\(XQ\), after choosing the soft-loss exponent below
\(\lambda_0\).
\end{proposition}

\begin{proof}
By \eqref{eq:tpc32-row-l1} and
\eqref{eq:tpc32-row-drift},
\[
 \sum_\alpha|\gamma_\alpha^{(i)}|\delta_H(m_\alpha)
 \ll X^\eps Q/L.
\]
Use this on the drift row, \eqref{eq:tpc32-row-l1} on the other row,
\eqref{eq:tpc32-prefix-envelope} for the ultra increment, and the
\(O_{\mathscr D}(J)\) orbit length.  Each of the two legs is
\[
 \ll X^\eps J(Q/L)Q
 \asymp X^\eps XQ/L.
\]
Renaming the harmless soft-loss exponent proves the claim.
\end{proof}

\begin{remark}[What is carried into the next sections]
The hard raw object is the single matched kernel
\eqref{eq:tpc32-exact-cutoff-shell}, not three unrelated absolute
sums.  The calibrated packet equals that object plus the rigorously
bounded drift contribution \eqref{eq:tpc32-drift-boundary}.  Every
factorization below retains the joint multiplier
\eqref{eq:tpc32-joint-multiplier}; boundedness alone is not treated as
a factorization theorem.
\end{remark}
